\section{Implementation}
\label{sec:implementation}

The toolkit is released as an open-source Python package
(\texttt{pip install robopsych}) with the source code committed to
\url{https://github.com/jrcruciani/robopsychology}. The architecture
is intentionally minimal so that the method can be re-implemented in
another language without surprises. This section describes the four
core components.

\subsection{CLI and prompt library}

The CLI (\texttt{robopsych} command, implemented in \texttt{cli.py})
exposes the diagnostic prompts and the orchestration glue. All prompts
live in YAML files under \texttt{src/robopsych/data/prompts/}, keyed
by identifier (\texttt{1.1}, \texttt{1.2}, ..., \texttt{4.3}). The
ratchet is run via \texttt{robopsych ratchet [--pure] [--behavioral]
[--coherence-judge MODEL]}.

Each diagnostic step is sent to the target model through a thin
provider abstraction (\texttt{providers.py}) that supports
Anthropic, OpenAI-compatible, and Google Gemini back ends. Provider
\texttt{send()} accepts optional \texttt{temperature} and
\texttt{response\_format} arguments and raises
\texttt{UnsupportedProviderOption} when a back end cannot honor a
request, letting the LLM-judge degrade gracefully (see
Section~\ref{sec:layer2}).

\subsection{Structured label parser}

The \texttt{labels.py} parser enforces the prompt grammar imposed by
rule~R2: every claim must be marked \texttt{[Observed]} or
\texttt{[Inferred]} and tagged with one of \texttt{[Model]},
\texttt{[Runtime]}, \texttt{[Conversation]}, \texttt{[Meta]}, or
\texttt{[Other]}. Unlabeled or malformed claims fall through a
fallback heuristic and are surfaced in the report so the practitioner
can spot dropped grammar before consuming downstream metrics.

\subsection{Two coherence analyzers}

The toolkit ships two coherence analyzers that share the
\texttt{CoherenceReport} contract so downstream report generation and
scoring code is analyzer-agnostic.

The \textbf{regex baseline} (\texttt{coherence.py}) counts surface
backward-reference patterns (``as I mentioned'', ``consistent with my
earlier'', ``building on my prior'') and surface contradiction
markers (``however I previously said'', ``upon reflection, I was
wrong'') over the response set. It is cheap, deterministic, and
useful as a floor.

The \textbf{LLM-judge analyzer} (\texttt{coherence\_llm.py}) implements
the four-layer pattern of \citet{llm_judge_skill}:

\begin{description}
    \item[\textbf{Layer~1} (deterministic pre-filter).] Splits each
        response into sentence-like units and drops fragments shorter
        than three words, questions, and \emph{hedged} statements
        (markers: \texttt{i think}, \texttt{maybe}, \texttt{perhaps},
        \texttt{i'm not sure}, \texttt{creo}, \texttt{tal vez},
        \texttt{quizás}, \texttt{puede que}). Hedged sentences are
        not commitments and must not be downgraded to contradictions
        by the judge.

    \item[\textbf{Layer~2} (judge discipline).]
        \label{sec:layer2}
        Calls the judge with \texttt{temperature=0} and requests
        \texttt{response\_format=\{"type": "json\_object"\}},
        falling back without the format hint if the provider raises
        \texttt{UnsupportedProviderOption}. The judge prompt embeds
        \emph{do-not-flag} examples (polite refinement, hedged
        statements, meaning-preserving reformulations) and grades
        contradictions by severity (high / medium / low).

    \item[\textbf{Layer~3} (automatic fallback).]
        \texttt{analyze\_coherence\_auto()} returns an
        \texttt{LLMCoherenceReport} with an \texttt{llm\_used} flag.
        When no judge is configured, the regex floor is used and the
        flag is set to \texttt{False}. The function accepts three
        documented \texttt{model\_config} shapes for testability and
        Azure parity.

    \item[\textbf{Layer~4} (exposed scoring weights).]
        \label{sec:layer4}
        Severity weights and the reference-credit / fresh-penalty
        knobs are exposed as a dictionary so calibration changes
        do not require re-prompting the judge.
\end{description}

\subsection{Reproducible case-study harness}

The repository's \texttt{validation/reproducible/} subtree hosts the
three case studies and a runner (\texttt{run\_case.py}) that exercises
the full pipeline against the live API. With \texttt{--runs N} the
runner replays each case $N$ times, writes per-run artifacts to
\texttt{artifacts/run-\{i\}/}, and aggregates them into
\texttt{distribution.json} (per-metric mean, standard deviation, min,
max, median; rate for boolean fields; frequency for list fields).
This is the infrastructure that drives the empirical claims in
Section~\ref{sec:validation}.

\todo{If room, add a small figure: architecture diagram with arrows
target $\to$ engine $\to$ ratchet $\to$ coherence analyzer $\to$
report. The drawing should fit on a third of a column. Source the
PDF in \texttt{figures/architecture.pdf}.}
